\documentclass{article}
\usepackage{graphicx}
\usepackage{fancyhdr}
\usepackage{geometry}
\usepackage{setspace}
\usepackage{tikz}
\usepackage[italian]{babel}
\usepackage[hidelinks]{hyperref}

% Margini della pagina
\geometry{a4paper, margin=1in}

% Intestazione personalizzata
\pagestyle{fancy}
\fancyhf{}
\fancyhead[L]{Code7Crusaders - Software Development Team}
\fancyhead[R]{\thepage}

% Spaziatura delle righe
\setstretch{1.2}

\begin{document}

% Sezione del titolo
\begin{titlepage}

    \AddToHookNext{shipout/background}{
    \begin{tikzpicture}[remember picture,overlay]
    \node at (current page.center) {
    %\includegraphics{background.png} %MODIFICA OVERLEAF
    \includegraphics{../../img/background.png}
    };
    \end{tikzpicture}
    }

    \centering
    \vspace*{2cm}

    %\includegraphics[width=0.3\textwidth]{7Crusaders_logo.png} % MODIFICA OVERWLOEAF
    \includegraphics[width=0.3\textwidth]{../../img/logo/7Crusaders_logo.png} % Aggiungi il logo qui
    \vspace{1cm}
    
    {\Huge \textbf{Code7Crusaders}}\\
    \vspace{0.5cm}
    {\Large Software Development Team}\\
    \vspace{2cm}
    
    {\large \textbf{Incontro del 18/10/2024 con Ergon}}\\
    \vspace{5cm}

    \textbf{Membri del Team:}\\
    Enrico Cotti Cottini, Gabriele Di Pietro, Tommaso Diviesti \\
    Francesco Lapenna, Matthew Pan, Eddy Pinarello, Filippo Rizzolo \\
    \vspace{0.5cm}
    
    \vspace{1cm}
\end{titlepage}

%Versioni
\begin{center}
    \textbf{Versioni}
    \\
    \begin{tabular}{|c|c|c|c|}
        \hline
        \textbf{Ver} & \textbf{Data} & \textbf{Autore} & \textbf{Descrizione}\\
        \hline
        1.0 & 28/10/2024 & Filippo Rizzolo & Approvazione documento \\ 
        \hline
        0.3 & 22/10/2024 & Filippo Rizzolo & Prima revisione del documento e controllo ortografia \\ 
        \hline
	    0.2 & 21/10/2024 & Enrico Cotti Cottini & Aggiunte Presenze \\
	    \hline
        0.1 & 18/10/2024 & Eddy Pinarello & Prima stesura del documento \\ 
        \hline
    \end{tabular}
\end{center}

% Indice
\newpage
\tableofcontents
\newpage

% Sezione Introduzione
\section{Registro Presenze}
\textbf{Piattaforma della riunione:} Zoom \\
\textbf{Ora di Inizio} 11:00\\
\textbf{Ora di Fine} 12:00\\
\\
\begin{tabular}{|c|c|c|}
    \hline
    \textbf{Componente} & \textbf{Ruolo} & \textbf{Presenza}\\
    \hline
    Enrico Cotti Cottini & Responsabile & Presente \\ 
    \hline
    Gabriele Di Pietro & Redattore & Presente \\ 
    \hline
    Tommaso Diviesti & Redattore & Presente \\ 
    \hline 
    Francesco Lapenna & Verificatore & Assente \\ 
    \hline
    Matthew Pan & Verificatore & Assente \\ 
    \hline 
    Eddy Pinarello & Redattore & Presente \\ 
    \hline 
    Filippo Rizzolo & Amministratore & Presente \\ 
    \hline 
\end{tabular}
\\
\newline
\newline
\begin{tabular}{|c|c|}
    \hline
    \textbf{Nome} & \textbf{Ruolo}\\
    \hline
    Gianluca Carlesso & Rappresentante Azienda \\
    \hline
    Anna Tieppo & Rappresentante Azienda \\
    \hline
\end{tabular}
\newpage

\section{Verbale dell'incontro}

\subsection{Utilizzo di .NET MAUI}
\textbf{Domanda:} L'utilizzo di .NET MAUI è un vincolo obbligatorio oppure ci sono altre alternative? \newline
\textit{Risposta:} No, è facoltativo. Utilizzate ciò che è più comodo per voi per ottenere il risultato.

\subsection{Fornitura di informazioni all'LLM}
\textbf{Domanda:} Come possiamo dare le informazioni all'LLM? Ci verrà fornita un'API o della documentazione? \newline
\textit{Risposta:} Noi vi forniamo un file contenente le informazioni del contesto di lavoro da cui l'AI apprenderà. Possiamo concordare il formato e saremo noi a recuperare le informazioni dai nostri clienti, che potrete utilizzare.

\subsection{Domande frequenti per l'allenamento dell'AI}
\textbf{Domanda:} Quali sono le domande più frequenti su cui l'AI deve essere allenata? \newline
\textit{Risposta:} A livello di domande, si può strutturare la chatbot in due modalità: una parte con domande predefinite, in cui l'AI deve comprendere e rispondere, e una seconda parte con l'AI generativa, dove l'LLM elabora la risposta basandosi solo sulla domanda.

\subsection{Chiarimento sull'AI}
\textbf{Domanda:} L'AI che useremo è già pre-addestrata? \newline
\textit{Risposta:} Sì, l'AI è già pre-addestrata. I modelli vengono allenati su grandi moli di dati. Nel nostro caso, la conoscenza fornita verrà trasmessa ai modelli, che analizzeranno ed elaboreranno le risposte.

\subsection{Utilizzo di modelli proprietari}
\textbf{Domanda:} Possiamo utilizzare modelli AI diversi da quelli consigliati, ad esempio modelli proprietari? \newline
\textit{Risposta:} Sì, potete utilizzare altri modelli. Vi proponiamo quelli perché sono progetti italiani e potrebbe essere interessante approfondirli. Se scegliete modelli proprietari, dovrete usare strumenti black box. Non c'è un vincolo specifico su queste tecnologie, ma noi consigliamo di usarle.

\subsection{Modello embedded a pagina 5}
\textbf{Domanda:} Come funziona il modello embedded descritto a pagina 5? \newline
\textit{Risposta:} Gli LLM utilizzano database vettoriali. Il testo fornito per l'apprendimento viene tokenizzato e rappresentato come un vettore, che poi viene salvato nel database vettoriale. Successivamente, l'LLM apprende da dati esterni che vengono forniti, a seconda del modello scelto. Gli API ricevono le domande dagli utenti, le tokenizzano e le passano all'LLM, che estrae le risposte dal database vettoriale. L'LLM utilizza un modello probabilistico per definire il contesto della domanda e fornire una risposta.

\subsection{Utilizzo di RAG come database vettoriale}
\textbf{Domanda:} Consigliate di usare RAG come database vettoriale? \newline
\textit{Risposta:} Sì, consigliamo RAG perché è uno dei più diffusi e ha molta documentazione disponibile. Posso eventualmente suggerirvi altri database.

\subsection{Rappresentazione dei dati in formato JSON}
\textbf{Domanda:} Nell'intermezzo tra database e LLM, è possibile rappresentare i dati sotto forma di JSON? \newline
\textit{Risposta:} Sì, questa è una possibilità che offre maggiore libertà. Avendo i dati in formato JSON, dovrete costruire l'interfacciamento e far comunicare i sistemi. Tuttavia, i dati forniti dall'azienda potrebbero non essere perfetti, quindi sarà necessaria una fase di pre-processing per eliminare eventuali dati anomali, dopo di che potete salvarli nel formato che preferite.

\subsection{Fasi di testing del prodotto}
\textbf{Domanda:} Nelle specifiche del progetto sono richieste anche delle fasi di testing del nostro prodotto. Ci sono delle linee guida o degli strumenti che ci consigliate? \newline
\textit{Risposta:} È preferibile fare i test durante lo sviluppo del progetto. Avrete domande con risposte attese e potrete creare test-set per verificare che il modello non presenti problemi di overfitting.

\subsection{Strumenti per colloqui e chiarimenti}
\textbf{Domanda:} Quale strumento usa l'azienda per eventuali colloqui e chiarimenti specifici, o ci si può incontrare dal vivo? \newline
\textit{Risposta:} Possiamo utilizzare Zoom per i colloqui o incontrarci in sede. Negli anni scorsi, i gruppi richiedevano incontri via Zoom e poi sono venuti anche in sede.

\subsection{Esperienza con i gruppi precedenti}
\textbf{Domanda:} Siccome non è il primo anno che proponete progetti ai nostri ex-colleghi, ci sono state difficoltà particolari per loro? \newline
\textit{Risposta:} Non hanno avuto grosse difficoltà. Sono sempre stati abbastanza autonomi e ci siamo sentiti settimanalmente via mail. Non ci sono state particolari difficoltà, anzi, siamo stati molto soddisfatti dal lavoro svolto.

\begin{table}[b]
	\begin{tabular}{@{}p{.5in}p{4in}@{}}
		Data:  & \hrulefill \\
			   &     		\\
			   &     		\\
		Firma: & \hrulefill \\
	\end{tabular}
	\end{table}

\end{document}